\thispagestyle{empty}
\begin{center}
{\small\textsc{Antern Studies \quad$\cdot$\quad Lecture Notes}}\\[14pt]
{\LARGE\bfseries Transformers from First Principles}\\[8pt]
{\large From one word to designing your own model}\\[12pt]
Ayush Singh \quad\textperiodcentered\quad Kirti Goyal\\[2pt]
{\small October 2026}
\end{center}
\vspace{4pt}

\section*{How to read these notes}

Why does a transformer look the way it does? You could memorise the diagram. I would rather you walk the road
the researchers walked and arrive at it yourself, one failure at a time. Every piece of the transformer
(attention, positions, the MLP, the residual stream, normalisation, the KV cache) exists because something
simpler broke. So in every chapter we first build the simple thing, watch it break \emph{with actual numbers},
and only then introduce the fix. If at the end of a section you feel ``I could have invented this myself'',
the section has done its job. Could you have invented attention in 2014, given what people knew then?
By Chapter 3 I think you will say yes.

\paragraph{Two readers at once.} These notes are written for someone who left school decades ago and for a
PhD student, at the same time. Every idea is climbed in three rungs:
\begin{itemize}
\item \textsc{Picture it}: an everyday picture you can hold in your head. No symbols.
\item \textsc{Say it exactly}: the same idea in mathematics, with every symbol defined and every shape written.
\item \textsc{Why it had to be this way}: the general principle underneath, the part that transfers to
problems that have nothing to do with language.
\end{itemize}
If a rung feels too easy, skip it. If it feels too hard, go one rung down and come back. You never need
anything you were not given on an earlier page.

\paragraph{How to use the boxes.} Each chapter opens with a \textbf{map} and a short \textbf{Before you
start} check of the maths it needs, with a refresher right there; do it, even if it looks familiar. When you
see \textbf{Predict first}, stop and write down a guess before reading on. A wrong guess is useful: it shows
you exactly which idea to look at again. \textbf{Where this picture breaks} tells you the limit of an
analogy, so the picture does not mislead you later. \textbf{A tempting wrong idea} names a mistake many
people make, says why it feels right, and shows why it is wrong. \textbf{Check yourself} questions have
their answers at the end of the chapter; answer first, then look. For a wrong option, the answer says which
wrong idea it came from. \textbf{Looking back} brings back questions from earlier chapters, because recalling
something after a gap is how it stays.

Why so many questions? Because a guess you wrote down, right or wrong, is what makes the answer stick.

\paragraph{Three rules hold throughout.}
\begin{enumerate}[label=(\roman*)]
\item \textbf{One running problem.} We start from a real message, \tok{waste of time}, and a real job:
reply to it well. The same tokens and the same tiny weights ($d=4$) follow us through every chapter, small
enough to compute with a pencil. Later the same steps are replayed on real models, GPT-2 and Qwen3-0.6B.
\item \textbf{Every number is real.} No number in these notes was typed by hand. Each one is produced by a
script in the companion repository and cross-checked (Appendix B lists them). Where we make a claim about the
world, such as ``counting word sequences runs out of data'', we measure it on a real corpus.
\item \textbf{Shapes on every equation.} Most confusion about transformers is shape confusion, not idea
confusion. Every equation carries its tensor shape in grey on the right.
\end{enumerate}

\paragraph{What you will be able to do.} By the end, without looking anything up, you should be able to
(a) run a full forward pass of a small transformer by hand, for a whole batch of sentences;
(b) explain, as a picture of arrows turning and stretching, what training does to every weight;
(c) read any model's \texttt{config.json} and compute its parameter count, training cost and KV-cache memory;
(d) design a transformer for a given compute, memory and latency budget, and defend every choice.

\paragraph{Contents.}
\begin{center}
\small
\begin{tabular}{@{}r@{\quad}l@{\qquad}r@{\quad}l@{}}
0 & Notation & 7 & From scores to words: temperature, sampling\\
1 & The problem, and 70 years of attempts & 8 & Learning: how the arrows get nudged\\
2 & Tokens, and words as points in space & 9 & Inference: the KV cache and its bill\\
3 & Attention: letting tokens talk & 10 & The toy scaled up: GPT-2, then Qwen3\\
4 & Order: from sinusoids to rotations (RoPE) & 11 & Designing your own transformer\\
5 & Many heads, the MLP, the residual stream & A & The calculus, checked numerically\\
6 & Normalisation, and why depth is hard & B & Reproducing every number\\
\end{tabular}
\end{center}
